% Source: https://artofproblemsolving.com/community/c3h1296592_optimization_in_squares
% Stable identity: aops-topic-1296592-q01
Put $A=(0,0)$, $B=(3,0)$, $P=(1,0)$, $Q=(2,0)$. For a point $E$ on a line not crossing segment $PQ$, $\angle PEQ$ is largest where the circle through $P$ and $Q$ is tangent to that line (points inside the circle see $PQ$ at a larger angle, points outside at a smaller one).
\begin{itemize}
\item Side $AD$ ($x=0$): the tangent circle has centre $(1.5,k)$ and radius $1.5$, with $0.25+k^2=2.25$, $k=\sqrt2$; it touches $AD$ at $E=(0,\sqrt2)$. By symmetry $BC$ gives $(3,\sqrt2)$.
\item Side $CD$ ($y=3$): the tangent circle has radius $3-k$ with $0.25+k^2=(3-k)^2$, so $k=\dfrac{35}{24}$ and radius $\dfrac{37}{24}>1.5$.
\end{itemize}
A chord of length $1$ subtends the inscribed angle $\theta$ with $\sin\theta=\dfrac{1}{2R}$, so the smaller circle ($R=1.5$) gives the larger angle. Hence $E=(0,\sqrt2)$ (or $(3,\sqrt2)$).

\begin{asy}
usepackage("amsmath");
size(4.5cm);
defaultpen(fontsize(10pt));
real s=3; pair A=(0,0),B=(3,0),C=(3,3),D=(0,3),P=(1,0),Q=(2,0),E=(0,sqrt(2));
draw(A--B--C--D--cycle,linewidth(0.9)); draw(circle((1.5,sqrt(2)),1.5),gray+dashed);
draw(P--E--Q,blue); dot(P);dot(Q);dot(E); label("$A$",A,SW); label("$B$",B,SE); label("$C$",C,NE); label("$D$",D,NW);
label("$P$",P,S); label("$Q$",Q,S); label("$E$",E,W);
label("\mbox{}",(min(currentpicture,user=true)+max(currentpicture,user=true))/2,black); // keeps the text after the figure black when LaTeX typesets the labels inline
\end{asy}

The area of $PEQ$ is $\dfrac{1}{2}\cdot PQ\cdot\sqrt2=\boxed{\dfrac{\sqrt2}{2}}$.
